\section{Retrieval and Capacity Proofs}

\subsection{Co-occurrence and the query schedule}
\label{sec:cooccurrence}
For every active record $u$, its target-rank interval contains exactly
the occupied targets whose proofs require $u$. Two records therefore
co-occur if and only if their intervals intersect. An intersection of
integer rank intervals contains an occupied rank. Laminarity turns this
condition into ancestry in the interval forest. Consequently a color
class is an antichain, and the records required by any target form a
chain. Conversely, every chain shares the nonempty interval of its
deepest member, so some target requires the entire chain. The maximum
proof width equals the forest height.

A widest proof needs a distinct color for each record. Coloring every
vertex by its forest depth attains that many colors in one traversal.
A proof has at most $h$ siblings. Its nonempty sibling subtrees are
disjoint and each contains an occupied leaf other than the target,
giving the additional bound $n-1$. To attain both bounds, put a target
at $0^h$, set $k=\min\{h,n-1\}$, and add the $k$ coordinates
$0^d1\,0^{h-d-1}$ for $0\le d<k$. If $n>h+1$, add any distinct unused
coordinates until the occupancy is $n$. The target already has $h$
nonempty siblings and retains them.

Proper coloring gives at most one required record in each database.
Backend correctness retrieves it; unused colors issue a dummy query.
The client retains only real answers and completes all other levels
with the public default chain. This reconstructs the usual membership
proof, including the standard root check against the pinned snapshot.

For the independent-instance privacy proof, condition on the setup leakage and fix two challenge
targets. In hybrid $j$, colors $1,\ldots,j$ use the indices for the
second target and the remaining colors use those for the first. Real
and dummy indices both pass through the same PIR interface. Index
privacy also covers a dummy mixture by averaging over its sampled
index. Independent client states and coins allow a reduction to
simulate the other colors. If $\epsilon_c$ bounds the distinguishing
advantage for color $c$, the hybrid argument gives
\begin{equation}
 \operatorname{Adv}_{\rm target}
 \le\sum_{c=1}^{m}\epsilon_c.
 \label{eq:supp-privacy-hybrid}
\end{equation}
For polynomial $m$ and negligible $\epsilon_c$, this is negligible.
The proof concerns the specified batch transcript and public setup;
it uses the model's independently held target value and matching root
and layout.

For a shared-state backend, the required assumption instead concerns
the complete batch-index vector, including shared evaluation keys and
fixed message lengths. Each target induces a distribution over valid
index vectors; joint query privacy makes the resulting server views
indistinguishable, also after averaging over dummy choices. Local
proof reconstruction adds no server-visible output. Thus the shared-state
realization inherits the backend's joint privacy bound, rather than
the independent-instance sum above. Both realizations inherit their
backend correctness-error bounds.

\subsection{Joint anchored-color capacity}
\label{sec:joint-anchor-proof}
Let $F$ be the record forest, with $N$ vertices and $m$ colors. Add a
virtual root $o$ at depth zero. Context $F_o$ is all of $F$; context
$F_x$ for a real vertex contains its strict descendants. Write
$N_x=|F_x|$ and $q_x=m-d(x)$. The colors on the root--$x$ chain are
unavailable throughout $F_x$, proving $N_x\le q_x U$ for every coloring
of capacity $U$. Define the context bound by
\begin{equation}
 B(F)=\max_{x:q_x>0}\left\lceil\frac{N_x}{q_x}\right\rceil.
 \label{eq:supp-context-bound}
\end{equation}
The virtual-root term is the average bound $\lceil N/m\rceil$.

The maximum union of $t$ antichains is a standard poset coverage
quantity~\cite{antichaincoverage}. Denote it here by $\beta_t$, with
$\beta_0=0$. Values add over forest components. For a rooted tree $T_v$,
consider whether its root belongs to the selected union:
\begin{equation}
 \beta_t(T_v)=\max\left\{\sum_w\beta_t(T_w),\;
  1+\sum_w\beta_{t-1}(T_w)\right\},\quad t\ge1.
 \label{eq:compact-beta}
\end{equation}
The sums range over children $w$ of $v$. Omitting $v$ leaves all $t$
colors available in every child subtree. Including it excludes its
color from all descendants. Both alternatives are attainable by
combining optimal child subsets on a common palette, which proves
the recurrence.

Let $P=(v_1,\ldots,v_k)$ be a contiguous path starting at a component
root of $F_x$, where $1\le k<q_x$. Its vertices anchor $k$ distinct
colors. Let $S_0$ contain the other context components, and let $S_i$
contain the subtrees branching from $v_i$ away from $v_{i+1}$.
All $k$ anchored colors are available in $S_0$; only the last $k-i$
are available in $S_i$. No anchored color is available below $v_k$.
The union of the anchored classes is therefore bounded by
\begin{equation}
 J_x(P)=k+\beta_k(S_0)+\sum_{i=1}^{k-1}\beta_{k-i}(S_i).
 \label{eq:supp-joint-union}
\end{equation}
This maximum is attained when individual color capacities are relaxed:
choose a maximizing subset in each side forest, color it with its
permitted suffix palette, and include the anchors. Side forests are
pairwise incomparable, and each suffix excludes exactly the anchors
that conflict with that region. Unselected records remain uncolored
in this auxiliary problem. The other $q_x-k$ classes of a complete
coloring contain at most $(q_x-k)U$ records. Hence
\begin{equation}
 U\ge\left\lceil\frac{N_x-J_x(P)}{q_x-k}\right\rceil.
 \label{eq:supp-joint-ratio}
\end{equation}
Taking the maximum over contexts and paths, together with context
averages, defines the joint certificate
\begin{equation}
 B_{\rm joint}(F)=\max\left\{B(F),\;
 \max_{x,P:1\le |P|<q_x}
 \left\lceil\frac{N_x-J_x(P)}{q_x-|P|}\right\rceil\right\}.
 \label{eq:supp-joint-definition}
\end{equation}
Here $U^*$ is the minimum maximum class size among proper $m$-colorings.
The preceding necessary inequalities prove $B_{\rm joint}\le U^*$.
The empty forest has bound zero, and empty maxima add no term.

For comparison, let $n_x$ count the occupied ranks of a context:
$n_o=n$, and $n_x=|I_{\mathsf{srv}}(x)|$ for a nonempty real context.
A class containing $v$ has at most
$a_{x,v}=1+n_x-|I_{\mathsf{srv}}(v)|$ records. Its other intervals
are nonempty, disjoint, and outside the interval of $v$.
Thus $J_x(P)\le\sum_{v\in P}a_{x,v}$. Define
\begin{equation}
 B_{\rm hier}=\max\left\{B,\;
 \max_{x,P:\,1\le |P|<q_x}
 \left\lceil\frac{N_x-\sum_{v\in P}a_{x,v}}{q_x-|P|}\right\rceil
 \right\}.
 \label{eq:supp-individual-hierarchy}
\end{equation}
It follows that $B\le B_{\rm hier}\le B_{\rm joint}$. Empty contexts
add no bound. A path with $k=q_x$ reaches total depth $m$; all its
side forests have height at most the available suffix size, so
$J_x(P)=N_x$. Such a path adds no positive-denominator test.

\subsection{Computing all joint tests in $O(Nm)$ time}
\label{sec:joint-evaluation}
The following evaluation also applies to a supplied palette
$m\ge\operatorname{height}(F)$. For each vertex, store its depth,
subtree size, and the child sums
\begin{equation}
 S_t(u)=\sum_{z\text{ child of }u}\beta_t(T_z),
 \qquad 0\le t\le m.
 \label{eq:supp-child-coverage}
\end{equation}
The virtual root needs only these sums, not a record value.
Equation~\eqref{eq:compact-beta} computes all values in postorder.
There are $N$ child edges including virtual-root edges, so this stage
costs $O(Nm)$ time and space.

Fix a terminal depth $d\in\{1,\ldots,m-1\}$. Set $Q_d(o)=0$ and, for
every edge $u\to w$ with $d(w)\le d$, compute
\begin{equation}
 Q_d(w)=Q_d(u)+1+S_{d-d(u)}(u)-\beta_{d-d(u)}(T_w).
 \label{eq:supp-joint-potential}
\end{equation}
The difference of child sums counts the side forest away from $w$.
For a terminal $v$ at depth $d$ and any strict ancestor context $x$,
telescoping gives
\begin{align}
 J_x(v)&=Q_d(v)-Q_d(x),\label{eq:supp-joint-telescope}\\
 q_x-|P|&=(m-d(x))-(d-d(x))=m-d.
 \label{eq:supp-joint-denominator}
\end{align}
Each edge contributes one anchor and the coverage of its off-path
forest by the remaining anchored suffix. This includes the other
context components on the first edge and excludes descendants of
the final anchor, exactly as in~\eqref{eq:supp-joint-union}.

Since the denominator is independent of $x$, all tests ending at $v$
reduce to
\begin{equation}
 \left\lceil
 \frac{\max_{x\prec v}\{N_x+Q_d(x)\}-Q_d(v)}{m-d}
 \right\rceil,
 \label{eq:supp-joint-scan}
\end{equation}
where $x\prec v$ ranges over strict ancestors, including $o$.
Maintain this ancestor maximum during the traversal, evaluating an
endpoint before inserting it into the maximum. Store an attaining
context with the maximum to recover a witness path. Each depth costs
$O(N)$ time and $O(N)$ temporary storage. Including preprocessing, the
total is $O(Nm)$ time and space, with
$O(\log(N+m+2))$-bit words. Context averages are retained explicitly;
the empty forest has bound zero. This computes a necessary certificate
without solving complete capacity feasibility.

\subsection{Exact assignment within an AB scan}
\label{sec:ab-sorting-assignment}
For one candidate subtree, let $s_i$ count its records in an allowed
color and $a_j$ count outside records in an allowed destination color.
Colors used by strict ancestors are excluded from this palette.
The assignment cost is $\sum_i(s_i+a_{\pi(i)})^2$.
For $s\le s'$ and $a\le a'$, the exchange difference is
\begin{equation}
 (s+a)^2+(s'+a')^2-(s+a')^2-(s'+a)^2
 =2(s'-s)(a'-a)\ge0.
 \label{eq:supp-count-exchange}
\end{equation}
Reverse pairing therefore minimizes the squared cost. It also
minimizes the full sorted load vector: the crossed pair has the same
sum and lies between the two original extreme loads. Every inversion
removal weakly decreases the descending sorted vector. Repeating
exchanges reaches reverse order from any assignment. If either
difference is zero, the two paired-load multisets agree.

Thus reverse sorting and a minimum-cost Hungarian assignment attain
the same minimum candidate load multiset. Equal-count labels can be
permuted differently, so their subsequent subtree assignments and
refinement trajectories need not coincide. A scan builds histograms
in $O(Nm)$ time and space, evaluates at most $N$ assignments using
$O(m\log(m+1))$ work each, and applies at most one selected move in
$O(N)$ time. At most $R_1$ scans cost $O(R_1Nm\log(m+1))$ time.
Strictly decreasing global sorted vectors ensure termination; a scan
without improvement establishes local optimality for these moves.
